\documentclass[uplatex,dvipdfmx,a4paper,11pt]{jsarticle}
\input{../../../templates/note.tex}
\noteprint{統計力学1　第1章「統計力学とは（熱力学の復習）」}{2026-11-25}
%% note-moto	problems\統計力学1\shou01\統計力学1-1.tex
\begin{document}
\begin{kirokuichiran}
  \ichirangyou{\hyperlink{bt-m-003-thermo}{1}}{L1・計算}{◎◎◎◎}{S}{2026-11-11 2回目}
  \ichirangyou{\hyperlink{bt-m-039-saikousei}{2}}{L1・計算}{◎◎◎◎}{S}{2026-11-11}
  \ichirangyou{\hyperlink{bt-m-056-maxwell}{3}}{L2・計算}{◎◎◎○}{C}{2026-11-11 ヒント2 (4)記述}
  \ichirangyou{\hyperlink{bt-m-057-legendre}{4}}{L2・計算}{○○◎◎}{C}{2026-11-11 ヒント2 (1)記述 (2)記述}
  \ichirangyou{5}{L3・導出}{}{未着手}{}
  \ichirangyou{6}{L4・導出}{}{未着手}{}
\end{kirokuichiran}

\youtenseiri
\begin{noteshikilist}
  \item 比熱からエネルギー（単位物質量あたり） % 出典: p.2
    \[ u(T)=\int_{0}^{T}c_{V}(T')\dd T'+u(0) \notetag{1-1} \]
  \item 比熱からエントロピー（単位物質量あたり） % 出典: p.2
    \[ s=\int_{0}^{T}\frac{c_{V}(T')}{T'}\dd T' \notetag{1-2} \]
  \item 比熱は自由エネルギーの2階微分（$n$ は物質量） % 出典: p.2
    \[ c_{V}=-\frac{T}{n}\left(\pdv[2]{F}{T}\right)_{V} \notetag{1-3} \]
\end{noteshikilist}

\mondaihen
\level{1}
% toubun: kiroku\toubun\統計力学1\003_2026-11-11.tex
\mondaibangou{1}
\begin{mondai}{003-thermo}{計算}{エントロピー $S(U,V)$ から温度・圧力・熱容量・自由エネルギーを求める熱力学関係式}
% 出典: 講義ノート p.1（1.1節），p.6 式(19)--(21)
ある系のエントロピーが内部エネルギー $U$ と体積 $V$ の関数として
\[S(U,V)=\frac{4}{3}\,a^{1/4}V^{1/4}U^{3/4}\]
で与えられている（$a$ は正の定数）。$1/T=(\partial S/\partial U)_{V}$ と $p/T=(\partial S/\partial V)_{U}$ を用いて，次の問いに答えよ。
\begin{shoumon}
  \item 内部エネルギー $U$ を $T$ と $V$ で表せ。
  \item 圧力 $p$ を $U$ と $V$ で表せ。
  \item 定積熱容量 $C_{V}$ を $T$ と $V$ で表せ。
  \item ヘルムホルツ自由エネルギー $F=U-TS$ を $T$ と $V$ で表せ。
\end{shoumon}
\notekiroku{2026-10-27 S ◎◎◎◎ → 2026-11-11 S ◎◎◎◎}
\begin{kotaeutsushi}{2026-11-11}
\toi{1}{◎}{}
\begin{align*}
\frac{1}{T}&=\frac{4}{3}a^{1/4}V^{1/4}\cdot\frac{3}{4}U^{-1/4}\\
&=\left(\frac{aV}{U}\right)^{1/4}
\end{align*}
\[ \therefore\ U=aVT^{4} \]
\toi{2}{◎}{}
\[ \frac{p}{T}=\frac{1}{3}a^{1/4}V^{-3/4}U^{3/4} \]
\begin{align*}
p&=\frac{1}{3}a^{1/4}V^{-3/4}U^{3/4}\cdot\left(\frac{U}{aV}\right)^{1/4}\\
&=\frac{1}{3}V^{-1}U^{1}=\frac{U}{3V}
\end{align*}
\toi{3}{◎}{}
\[ C_{V}=\left(\pdv{U}{T}\right)_{V}=4aVT^{3} \]
\toi{4}{◎}{}
\begin{align*}
F&=aVT^{4}-\left(\frac{U}{aV}\right)^{1/4}\cdot\frac{4}{3}a^{1/4}V^{1/4}U^{3/4}\\
&=aVT^{4}-\frac{4}{3}aVT^{4}\\
&=-\frac{1}{3}aVT^{4}
\end{align*}
\end{kotaeutsushi}
\end{mondai}

% toubun: kiroku\toubun\統計力学1\039_2026-11-11.tex
\mondaibangou{2}
\begin{mondai}{039-saikousei}{計算}{比熱を積分して内部エネルギー・エントロピー・ヘルムホルツ自由エネルギーを組み立て，比熱が自由エネルギーの2階微分であることを確かめる}
ある固体の体積を一定に保って比熱を測ると，低温で \qty{1}{\mole} あたりの定積比熱が
  \[ c_{V}(T)=\gamma T+aT^{3} \]
  （$\gamma$，$a$ は正の定数）と表された。\qty{0}{\kelvin} での \qty{1}{\mole} あたりの内部エネルギーを $u(0)$ とし，エントロピーは $s(0)=0$ とする。この温度範囲で上の式が成り立つとして，次の問いに答えよ。
  \begin{shoumon}
    \item \qty{1}{\mole} あたりの内部エネルギー $u(T)$ を求めよ。
    \item \qty{1}{\mole} あたりのエントロピー $s(T)$ を求めよ。
    \item \qty{1}{\mole} あたりのヘルムホルツ自由エネルギー $f(T):=u-Ts$ を求めよ。
    \item (3) の $f(T)$ を体積一定で $T$ について2回微分し，
      \[ -T\pdv[2]{f}{T} \]
      が $c_{V}(T)$ に一致することを確かめよ（\noteshiki{1-3} を \qty{1}{\mole} あたりに書いたもの）。
  \end{shoumon}
\notekiroku{2026-11-11 S ◎◎◎◎}
\begin{kotaeutsushi}{2026-11-11}
\toi{1}{◎}{}
\[ \left(\pdv{u}{T}\right)_{V}=\gamma T+aT^{3} \]
\[ \therefore\ u(T)=\frac{1}{2}\gamma T^{2}+\frac{1}{4}aT^{4}+u(0) \]
\toi{2}{◎}{}
\begin{align*}
s(T)&=\int_{0}^{T}\frac{c_{V}(T')}{T'}\dd T'\\
&=\int_{0}^{T}\kakimachigai{\gamma+aT'^{2}}{(\gamma+aT'^{2})}\dd T'\\
&=\gamma T+\frac{1}{3}aT^{3}
\end{align*}
\toi{3}{◎}{}
\begin{align*}
f(T)&=\frac{1}{2}\gamma T^{2}+\frac{1}{4}aT^{4}+u(0)\\
&\qquad-\gamma T^{2}-\frac{1}{3}aT^{4}\\
&=u(0)-\frac{1}{2}\gamma T^{2}-\frac{1}{12}aT^{4}
\end{align*}
\toi{4}{◎}{}
\begin{align*}
-T\pdv[2]{f}{T}&=-T\pdv{}{T}\left(-\gamma T-\frac{1}{3}aT^{3}\right)\\
&=-T(-\gamma-aT^{2})\\
&=\gamma T+aT^{3}\\
&=c_{V}(T)
\end{align*}
一致。
\end{kotaeutsushi}
\end{mondai}

\level{2}
% toubun: kiroku\toubun\統計力学1\056_2026-11-11.tex
\mondaibangou{3}
\begin{mondai}{056-maxwell}{計算}{マクスウェルの関係式による $(\partial U/\partial V)_{T}$ の計算，状態方程式と熱容量からの内部エネルギー・エントロピーの組み立て，自由膨張での温度変化とエントロピーの増加}
物質量 $n$ の気体の状態方程式が，ファン・デル・ワールスの式
\[ p=\frac{nRT}{V-nb}-\frac{an^{2}}{V^{2}} \]
（$a$，$b$ は正の定数）で与えられるとする。体積が十分大きい極限ではこの気体は単原子の理想気体に戻り，定積熱容量は $3nR/2$ になる。
\begin{shoumon}
  \item ギブスの関係式とマクスウェルの関係式を用いて
  \[ \left(\pdv{U}{V}\right)_{T}=T\left(\pdv{p}{T}\right)_{V}-p \]
  を示し，この気体について右辺を計算せよ。
  \item $(\partial C_{V}/\partial V)_{T}$ を $p$ の温度についての2階微分で表し，この気体では $C_{V}=3nR/2$ が任意の $V$ で成り立つことを示せ。さらに，$V\to\infty$ で $U\to3nRT/2$ となるように内部エネルギー $U(T,V)$ を求めよ。
  \item エントロピー $S(T,V)$ を求めよ（加える定数は $S_{0}$ と書いてよい）。
  \item Ar 気体 $\qty{1.00}{\mole}$（$a=\qty{0.1355}{\pascal\metre\tothe{6}\per\mole\squared}$，$b=\qty{3.20e-5}{\cubic\metre\per\mole}$）が，温度 \qty{300}{\kelvin} で体積 $V=\qty{1.00e-3}{\cubic\metre}$ の断熱容器の中にあり，真空の部屋との仕切りを取り除くと全体積 $2V$ に広がった（自由膨張）。膨張後の温度の変化 $\Delta T$ とエントロピーの変化 $\Delta S$ を有効数字3桁で求めよ。$R=\qty{8.314}{\joule\per\kelvin\per\mole}$ とする。
\end{shoumon}
\notekiroku{2026-11-11 C ◎◎◎○ (ヒント2)}
\begin{kotaeutsushi}{2026-11-11}
\toi{1}{◎}{}
（図：熱力学の四角形。上の辺に $S$，$U$，$V$，左右に $H$，$F$，下の辺に $p$，$G$，$T$ を置き，対角線に矢印）
\begin{align*}
\dd U&=\pdv{U}{S}\dd S+\pdv{U}{V}\dd V\\
&=T\,\dd S-p\,\dd V
\end{align*}
\[ \left(\pdv{U}{V}\right)_{T}=T\left(\pdv{S}{V}\right)_{T}-p \]
\[ \kakimachigai{\left(\pdv{S}{U}\right)_{T}}{\left(\pdv{S}{V}\right)_{T}}=\left(\pdv{p}{T}\right)_{V} \]
\[ \left(\pdv{U}{V}\right)_{T}=T\left(\pdv{p}{T}\right)_{V}-p \]
\begin{align*}
\left(\pdv{U}{V}\right)_{T}&=T\left(\frac{nR}{V-nb}\right)\\
&\qquad-\frac{nRT}{V-nb}+\frac{an^{2}}{V^{2}}\\
&=\frac{an^{2}}{V^{2}}
\end{align*}
\toi{2}{◎}{}
\[ C_{V}=\left(\pdv{U}{T}\right)_{V} \]
\begin{align*}
\left(\pdv{C_{V}}{V}\right)_{T}&=\pdv{}{T}\left(\left(\pdv{U}{V}\right)_{T}\right)_{V}\\
&=\pdv{}{T}\left(T\left(\pdv{p}{T}\right)_{V}-p\right)_{V}\\
&=\kakimachigai{\pdv{}{T}\cdot T}{\pdv{T}{T}}\left(\pdv{p}{T}\right)_{V}\\
&\qquad+T\left(\pdv[2]{p}{T}\right)_{V}-\left(\pdv{p}{T}\right)_{V}\\
&=T\left(\pdv[2]{p}{T}\right)_{V}
\end{align*}
\[ T\left(\pdv[2]{p}{T}\right)=0\qquad\therefore\ \left(\pdv{C_{V}}{V}\right)_{T}=0 \]
$C_{V}=3nR/2$ が任意の $V$ で成り立つ。
\[ \left(\pdv{U}{V}\right)_{T}=\frac{an^{2}}{V^{2}} \]
\[ \therefore\ U(T,V)=-\frac{an^{2}}{V}+U(T) \]
$V\to\infty$ で $U\to3nRT/2$ より $U(T)=3nRT/2$。
\[ U(T,V)=\frac{3nRT}{2}-\frac{an^{2}}{V} \]
\toi{3}{◎}{}
（最初の2つの試みは×で消して，次のように書き直した）
\[ \dd S=\left(\pdv{S}{T}\right)_{V}\dd T+\left(\pdv{S}{V}\right)_{T}\dd V \]
\[ \dd U=T\,\dd S-p\,\dd V \]
（脇で計算）
\[ \left(\pdv{S}{V}\right)_{T}=\left(\pdv{p}{T}\right)_{V}=\frac{nR}{V-nb} \]
\[ \dd S=\frac{1}{T}\dd U+\frac{p}{T}\dd V \]
\[ \left(\pdv{S}{T}\right)_{V}=\frac{1}{T}\left(\pdv{U}{T}\right)_{V}+\frac{p}{T}\cdot0 \]
\[ \dd S=\frac{C_{V}}{T}\dd T+\frac{nR}{V-nb}\dd V \]
\begin{align*}
S(T,V)&=\frac{3nR}{2}\int\frac{\dd T}{T}+\int\frac{nR}{V-nb}\dd V\\
&=\frac{3nR}{2}\log T+nR\log(V-nb)+S_{0}
\end{align*}
\toi{4}{○}{記述}
\[ \nuke{\text{熱も仕事も }0}\quad\Delta U=0 \]
\begin{align*}
U(T_{1},V)&=\frac{3nRT_{1}}{2}-\frac{an^{2}}{V}\\
U(T_{2},2V)&=\frac{3nRT_{2}}{2}-\frac{an^{2}}{2V}
\end{align*}
\[ \therefore\ \frac{3nR(T_{2}-T_{1})}{2}-\frac{an^{2}}{2V}+\frac{an^{2}}{V}=0 \]
\begin{align*}
3nR\Delta T&=-\frac{an^{2}}{V}\\
\Delta T&=-\frac{an}{3RV}\\
&=\qty{-5.43}{\kelvin}
\end{align*}
\nuke{$S$ は状態量}
\begin{align*}
\Delta S&=S(T_{2},2V)-S(T_{1},V)\\
&=\frac{3nR}{2}\log\frac{T_{2}}{T_{1}}+nR\log\frac{2V-nb}{V-nb}\\
&=\qty{5.67}{\joule\per\kelvin}
\end{align*}
\naze{$\Delta U=0$ の理由（断熱容器なので熱が 0，真空へ広がるので外への仕事も 0）と，自由膨張は不可逆過程だが $S$ は状態量なので，同じ始状態と終状態を結ぶ準静的な経路で求めた (3) の式の差で $\Delta S$ を計算してよいこと（答案に必要な論点）が，式にも文にも無い。式と数値は正しい。}
\end{kotaeutsushi}
\begin{kizuki}
\item \mae{第1章（旧版）問題2 (4)・10-27} 不可逆過程の変化量は，結果を出す式の前に「なぜその式を使ってよいか」を一言書く。自由膨張では熱も仕事も 0 なので $U$ が保存し，$S$ は状態量なので準静的な経路で求めた (3) の式の差で $\Delta S$ を求めてよい。前回も $2^{N}$ になる理由を書かなかったので，理由の1行を式より先に置く。
\item 自由膨張で保存するのは $U$ であって $T$ ではない。$(\partial U/\partial V)_{T}=an^{2}/V^{2}>0$ なので，$U$ 一定のまま体積が増えると温度が下がる。自由膨張を準静的断熱過程と取り違えて $\Delta S=0$ としない。
\end{kizuki}
\begin{tadashii}
\houshin ギブスの関係式から $(\partial U/\partial V)_{T}$ を $(\partial S/\partial V)_{T}$ で書き，マクスウェルの関係式で状態方程式だけで計算できる形にする。$C_{V}$ の体積依存性も状態方程式から決まるので，体積の大きい極限の値と合わせて $U$ と $S$ を組み立てる。自由膨張は $U$ 一定の不可逆過程なので，前後の状態を $U$ の保存で結び，$S$ は状態量として差をとる。
\keisan \shoumonkai{1} $\dd U=T\,\dd S-p\,\dd V$ で $T$ を一定にして $\dd V$ で割ると
\[ \left(\pdv{U}{V}\right)_{T}=T\left(\pdv{S}{V}\right)_{T}-p \]
$\dd F=-S\,\dd T-p\,\dd V$ より $S=-(\partial F/\partial T)_{V}$，$p=-(\partial F/\partial V)_{T}$ で，2階微分の順序を入れ替えると $(\partial S/\partial V)_{T}=(\partial p/\partial T)_{V}$。よって示すべき式が得られる。この気体では $(\partial p/\partial T)_{V}=nR/(V-nb)$ なので
\[ \left(\pdv{U}{V}\right)_{T}=\frac{nRT}{V-nb}-\left(\frac{nRT}{V-nb}-\frac{an^{2}}{V^{2}}\right)=\frac{an^{2}}{V^{2}} \]
\shoumonkai{2} (1) の式を $V$ 一定で $T$ について微分し，偏微分の順序を入れ替える。
\[ \left(\pdv{C_{V}}{V}\right)_{T}=\pdv{}{T}\left[T\left(\pdv{p}{T}\right)_{V}-p\right]_{V}=\left(\pdv{p}{T}\right)_{V}+T\left(\pdv[2]{p}{T}\right)_{V}-\left(\pdv{p}{T}\right)_{V}=T\left(\pdv[2]{p}{T}\right)_{V} \]
$p$ は $T$ の1次式なのでこれは $0$ で，$C_{V}$ は $V$ によらない。$V\to\infty$ での値が $3nR/2$（$T$ によらない）なので，任意の $V$ で $C_{V}=3nR/2$。$U$ は $(\partial U/\partial T)_{V}=3nR/2$ と (1) を満たすので
\[ U(T,V)=\frac{3}{2}nRT-\frac{an^{2}}{V} \]
（$V\to\infty$ で $3nRT/2$ になるので，加える定数は $0$）。
\shoumonkai{3} $\dd S=(\partial S/\partial T)_{V}\dd T+(\partial S/\partial V)_{T}\dd V$ で，$(\partial S/\partial T)_{V}=C_{V}/T$，$(\partial S/\partial V)_{T}=(\partial p/\partial T)_{V}=nR/(V-nb)$。積分して
\[ S(T,V)=\frac{3}{2}nR\log T+nR\log(V-nb)+S_{0} \]
\shoumonkai{4} 自由膨張では外部との熱も仕事も $0$ なので $U$ が保存する。膨張前の温度を $T_{1}=\qty{300}{\kelvin}$，膨張後を $T_{2}$ として
\[ \frac{3}{2}nRT_{1}-\frac{an^{2}}{V}=\frac{3}{2}nRT_{2}-\frac{an^{2}}{2V} \]
より
\[ \Delta T=T_{2}-T_{1}=-\frac{an}{3RV}=-\frac{\num{0.1355}\times\num{1.00}}{3\times\num{8.314}\times\num{1.00e-3}}\,\unit{\kelvin}=\qty{-5.43}{\kelvin} \]
$T_{2}=\qty{294.57}{\kelvin}$ である。エントロピーは状態量なので (3) の式の差をとる。
\[ \Delta S=\frac{3}{2}nR\log\frac{T_{2}}{T_{1}}+nR\log\frac{2V-nb}{V-nb} \]
第1項は $\num{12.47}\times\log(\num{0.98189})\,\unit{\joule\per\kelvin}=\qty{-0.228}{\joule\per\kelvin}$，第2項は
\[ \num{8.314}\times\log\frac{\num{1.968e-3}}{\num{0.968e-3}}\,\unit{\joule\per\kelvin}=\num{8.314}\times\num{0.7095}\,\unit{\joule\per\kelvin}=\qty{5.899}{\joule\per\kelvin} \]
よって $\Delta S=\qty{5.67}{\joule\per\kelvin}$。
\kotae{\shoumonkai{1} $(\partial U/\partial V)_{T}=an^{2}/V^{2}$ \shoumonkai{2} $(\partial C_{V}/\partial V)_{T}=T(\partial^{2}p/\partial T^{2})_{V}=0$，$U=3nRT/2-an^{2}/V$ \shoumonkai{3} $S=\frac{3}{2}nR\log T+nR\log(V-nb)+S_{0}$ \shoumonkai{4} $\Delta T=\qty{-5.43}{\kelvin}$，$\Delta S=\qty{5.67}{\joule\per\kelvin}$}
\begin{ronten}
  \item \shoumonkai{1} $(\partial S/\partial V)_{T}=(\partial p/\partial T)_{V}$ が $F$ の2階微分の順序の入れ替えから出ることを書く。
  \item \shoumonkai{2} $C_{V}$ が $V$ によらないことを示したうえで，体積の大きい極限の値を使って $C_{V}$ と $U$ の定数を決めることを述べる。
  \item \shoumonkai{4} 自由膨張は不可逆過程だが，$S$ は状態量なので，同じ始状態と終状態を結ぶ準静的な経路の式 (3) で差を計算してよいことを述べる。
\end{ronten}
\kenzan 次元：$an^{2}/V^{2}$ は圧力の次元で，$(\partial U/\partial V)_{T}$（エネルギー/体積）と一致する。$an/(RV)$ は $(\unit{\pascal\cubic\metre\per\mole})/(\unit{\joule\per\kelvin\per\mole})=\unit{\kelvin}$。特殊ケース：$a=b=0$ では $\Delta T=0$，$\Delta S=nR\log2=\qty{5.76}{\joule\per\kelvin}$ で理想気体に戻る。$\Delta S>0$ は孤立系の不可逆過程でエントロピーが増えることと整合する。
\tsumazuki $(\partial U/\partial V)_{T}$ を $-p$ としてしまう誤りが多い（それは $(\partial U/\partial V)_{S}$）。自由膨張を準静的断熱過程と取り違えて $\Delta S=0$ としない。
\end{tadashii}
\end{mondai}

% toubun: kiroku\toubun\統計力学1\057_2026-11-11.tex
\mondaibangou{4}
\begin{mondai}{057-legendre}{計算}{基本関係式 $S(U,V,N)$ から示強変数・状態方程式を求め，ルジャンドル変換で $F$・$G$ を作ること，オイラーの関係式とギブス--デュエムの関係式}
ある仮想的な物質の基本関係式が
\[ S(U,V,N)=A\,(NVU)^{1/3} \]
（$A$ は正の定数）で与えられるとする。
\begin{shoumon}
  \item $S$ が示量的であることを確かめよ。また，温度 $T$・圧力 $p$・化学ポテンシャル $\mu$ を $U$，$V$，$N$，$S$ で表し，オイラーの関係式 $U=TS-pV+\mu N$ が成り立つことを確かめよ。
  \item 内部エネルギー $U(T,V,N)$，状態方程式 $p(T,V,N)$，ヘルムホルツ自由エネルギー $F(T,V,N)$ を求めよ。
  \item ギブス自由エネルギー $G(T,p,N)$ を求め，$G=\mu N$ を確かめよ。
  \item (3) の $\mu$ を $T$ と $p$ の関数と見て，ギブス--デュエムの関係式から導かれる $(\partial\mu/\partial T)_{p}=-S/N$ と $(\partial\mu/\partial p)_{T}=V/N$ が成り立つことを確かめよ。
\end{shoumon}
\notekiroku{2026-11-11 C ○○◎◎ (ヒント2)}
\begin{kotaeutsushi}{2026-11-11}
\toi{1}{○}{記述}
\begin{align*}
\lambda S&=S(\lambda U,\lambda V,\lambda N)\\
&=A(\lambda^{3}NVU)^{1/3}\\
&=\lambda A(NVU)^{1/3}
\end{align*}
より示量的。
\[ \dd U=T\,\dd S-p\,\dd V+\mu\,\dd N \]
積分すると
\[ U=TS-pV+\mu N \]
（続き）
\[ \dd U=T\,\dd S-p\,\dd V+\mu\,\dd N \]
\[ \therefore\ \dd S=\frac{1}{T}\dd U+\frac{p}{T}\dd V-\frac{\mu}{T}\dd N \]
\begin{gather*}
\left(\pdv{S}{U}\right)_{V}=\frac{1}{T},\quad\left(\pdv{S}{V}\right)_{U}=\frac{p}{T},\\
\left(\pdv{S}{N}\right)_{U,V}=-\frac{\mu}{T}
\end{gather*}
\begin{align*}
\left(\pdv{S}{U}\right)_{V}&=A(NV)^{1/3}\cdot\frac{1}{3}U^{-2/3}=\frac{S}{3U}\\
\left(\pdv{S}{V}\right)_{U}&=\frac{S}{3V},\quad\left(\pdv{S}{N}\right)_{U,V}=\frac{S}{3N}
\end{align*}
\[ \therefore\ T=\frac{3U}{S},\quad p=\frac{S}{3V}\cdot\frac{3U}{S}=\frac{U}{V} \]
\[ \therefore\ \mu=-\frac{S}{3N}\cdot\frac{3U}{S}=-\frac{U}{N} \]
\nuke{$TS-pV+\mu N=3U-U-U=U$}
\naze{オイラーの関係式は $U=TS-pV+\mu N$ と書いただけで，求めた $T$・$p$・$\mu$ を代入していない。「確かめよ」なので，$TS-pV+\mu N=3U-U-U=U$ の一行で両辺が一致することまで書く。}
\toi{2}{○}{記述}
\[ \frac{1}{T}=\left(\pdv{S}{U}\right)_{V}=\frac{1}{3}A(NV)^{1/3}U^{-2/3} \]
\[ \frac{U^{2/3}}{T}=\frac{1}{3}A(NV)^{1/3} \]
\[ U(T,V,N)=\left(\frac{1}{3}AT\right)^{3/2}(NV)^{1/2} \]
\begin{align*}
\ayamari{\left(\pdv{p}{T}\right)_{V}}{\frac{p}{T}}&=\ayamari{\left(\pdv{S}{V}\right)_{T}}{\left(\pdv{S}{V}\right)_{U}}\\
&=\frac{1}{3}A(NU)^{1/3}V^{-2/3}
\end{align*}
\naze{右辺は $S(U,V,N)$ を $U$ 一定で $V$ について微分したもので，$(\partial S/\partial V)_{U}=p/T$ である。マクスウェルの関係式の $(\partial S/\partial V)_{T}=(\partial p/\partial T)_{V}$ は $T$ 一定の微分で別の量（この問題では $(\partial p/\partial T)_{V}=\frac{3}{2}p/T$）。次の行は $p=T\cdot(p/T)$ として正しく，$U$・$p$・$F$ の答えは正しい。}
\begin{align*}
p(T,V,N)&=\frac{1}{3}AT\,N^{1/3}V^{-2/3}U^{1/3}\\
&=\frac{1}{3}AT\,N^{1/3}V^{-2/3}\\
&\qquad\times\left(\frac{1}{3}AT(NV)^{1/3}\right)^{1/2}\\
&=\left(\frac{1}{3}AT\right)^{3/2}\left(\frac{N}{V}\right)^{1/2}
\end{align*}
\begin{align*}
F(T,V,N)&=U-TS\\
&=\left(\frac{1}{3}AT\right)^{3/2}(NV)^{1/2}-AT(NVU)^{1/3}\\
&=\left(\frac{1}{3}AT\right)^{3/2}(NV)^{1/2}\\
&\qquad-3\left(\frac{1}{3}AT\right)^{3/2}(NV)^{1/2}\\
&=-2\left(\frac{1}{3}AT\right)^{3/2}(NV)^{1/2}
\end{align*}
\toi{3}{◎}{}
\[ G=F+pV \]
\[ G(T,p,N)=-2\left(\frac{1}{3}AT\right)^{3/2}(NV)^{1/2}+pV \]
\[ p=\left(\frac{1}{3}AT\right)^{3/2}\left(\frac{N}{V}\right)^{1/2} \]
\[ p^{2}=\left(\frac{1}{3}AT\right)^{3}\frac{N}{V} \]
\[ V=\left(\frac{1}{3}AT\right)^{3}\cdot\frac{N}{p^{2}} \]
\begin{align*}
G(T,p,N)&=-2\left(\frac{1}{3}AT\right)^{3/2}\left(\frac{1}{3}AT\right)^{3/2}\\
&\qquad\times\left(\frac{N^{2}}{p^{2}}\right)^{1/2}+\left(\frac{1}{3}AT\right)^{3}\cdot\frac{N}{p}\\
&=-\left(\frac{1}{3}AT\right)^{3}\frac{N}{p}\\
&=\underbrace{-\left(\frac{1}{3}AT\right)^{3}p^{-1}}_{\mu}\cdot N
\end{align*}
\[ \mu=-\frac{U}{N}=-\left(\frac{1}{3}AT\right)^{3/2}\left(\frac{V}{N}\right)^{1/2} \]
\begin{align*}
\mu&=-\left(\frac{1}{3}AT\right)^{3}p^{-1}\\
&=-\left(\frac{1}{3}AT\right)^{3}\cdot\left(\frac{1}{3}AT\right)^{-3/2}\left(\frac{N}{V}\right)^{-1/2}\\
&=-\left(\frac{1}{3}AT\right)^{3/2}\left(\frac{V}{N}\right)^{1/2}
\end{align*}
で一致。
\toi{4}{◎}{}
\[ \mu=-\left(\frac{1}{3}AT\right)^{3}p^{-1} \]
\begin{align*}
\left(\pdv{\mu}{T}\right)_{p}&=\kakimachigai{-AT^{2}\cdot p^{-1}}{-\frac{A^{3}}{9}T^{2}p^{-1}}\\
&=-\frac{3}{T}\cdot\frac{1}{p}\left(\frac{1}{3}AT\right)^{3}
\end{align*}
\[ \left(\pdv{\mu}{p}\right)_{T}=+\left(\frac{1}{3}AT\right)^{3}p^{-2} \]
\begin{align*}
-\frac{S}{N}&=-\frac{3U}{NT}=\frac{3\mu}{T}\\
&=-\frac{3}{T}\cdot\frac{1}{p}\left(\frac{1}{3}AT\right)^{3}
\end{align*}
より一致。
\begin{align*}
\frac{V}{N}&=\frac{U}{Np}=-\frac{\mu}{p}\\
&=+\left(\frac{1}{3}AT\right)^{3}p^{-2}
\end{align*}
より一致。
\end{kotaeutsushi}
\begin{kizuki}
\item 基本関係式 $S(U,V,N)$ の偏微分は残りの変数を固定したもので，$(\partial S/\partial V)_{U,N}=p/T$ である。マクスウェルの関係式 $(\partial S/\partial V)_{T}=(\partial p/\partial T)_{V}$ は $\dd F=-S\,\dd T-p\,\dd V$ から出る $T$ 一定の微分で，別の量である。$(\partial p/\partial T)_{V}$ に $T$ を掛けて $p$ になるのは $p$ が $T$ の1次式のときだけで，この問題の $p\propto T^{3/2}$ では成り立たない。
\item \mae{第1章（旧版）問題4 (3)(4)・10-27} 記号の書き違いである。前回は $S'$ のプライムを落とし，今回は固定する変数を取り違えた（問題3 (1) でも $(\partial S/\partial V)_{T}$ を $(\partial S/\partial U)_{T}$ と書いた）。偏微分を書いたら，固定する変数がその式の出どころ（どの関数を何で微分したか）と合っているかを見直す。
\item 「確かめよ」の問いは，求めた $T$・$p$・$\mu$ を代入して $TS-pV+\mu N=3U-U-U=U$ のように両辺が一致する一行まで書く。
\end{kizuki}
\begin{tadashii}
\houshin 基本関係式の偏微分で示強変数を求め，$1/T$ の式を解いて $U(T,V,N)$ にする。$F=U-TS$，$G=F+pV$ を自然な変数で書き直し，オイラーの関係式・ギブス--デュエムの関係式を具体的に確かめる。
\keisan \shoumonkai{1} $S(\lambda U,\lambda V,\lambda N)=A(\lambda^{3}NVU)^{1/3}=\lambda S(U,V,N)$ なので示量的である。偏微分は
\[ \frac{1}{T}=\pdv{S}{U}=\frac{S}{3U},\qquad \frac{p}{T}=\pdv{S}{V}=\frac{S}{3V},\qquad -\frac{\mu}{T}=\pdv{S}{N}=\frac{S}{3N} \]
よって $T=3U/S$，$p=TS/(3V)=U/V$，$\mu=-TS/(3N)=-U/N$。これらを代入すると
\[ TS-pV+\mu N=3U-U-U=U \]
でオイラーの関係式が成り立つ。
\shoumonkai{2} $1/T=(A/3)(NV)^{1/3}U^{-2/3}$ を $U$ について解くと
\[ U=\left(\frac{AT}{3}\right)^{3/2}(NV)^{1/2} \]
$p=U/V$ より
\[ p=\left(\frac{AT}{3}\right)^{3/2}\left(\frac{N}{V}\right)^{1/2} \]
$TS=3U$ なので
\[ F=U-TS=-2U=-2\left(\frac{AT}{3}\right)^{3/2}(NV)^{1/2} \]
（確認：$-(\partial F/\partial V)_{T,N}=(AT/3)^{3/2}N^{1/2}V^{-1/2}=p$。）
\shoumonkai{3} 状態方程式を $V$ について解くと $V=N(AT/3)^{3}/p^{2}$。$pV=U$ なので $G=F+pV=-2U+U=-U$ で，$U=pV=N(AT/3)^{3}/p$ を使うと
\[ G(T,p,N)=-\frac{N}{p}\left(\frac{AT}{3}\right)^{3} \]
(1) の $\mu=-U/N$ に $U=N(AT/3)^{3}/p$ を入れると $\mu=-(AT/3)^{3}/p$ で，$G=\mu N$ が成り立つ。
\shoumonkai{4} $\mu(T,p)=-(A/3)^{3}T^{3}/p$ を微分する。
\[ \left(\pdv{\mu}{T}\right)_{p}=-\frac{3}{T}\cdot\frac{1}{p}\left(\frac{AT}{3}\right)^{3},\qquad \left(\pdv{\mu}{p}\right)_{T}=\frac{1}{p^{2}}\left(\frac{AT}{3}\right)^{3} \]
一方，$S/N=3U/(NT)=(3/T)(AT/3)^{3}/p$，$V/N=(AT/3)^{3}/p^{2}$ なので，$(\partial\mu/\partial T)_{p}=-S/N$，$(\partial\mu/\partial p)_{T}=V/N$ が成り立つ。これは $N\,\dd\mu=-S\,\dd T+V\,\dd p$（ギブス--デュエムの関係式）を $N$ で割ったものである。
\kotae{\shoumonkai{1} $T=3U/S$，$p=U/V$，$\mu=-U/N$（$TS-pV+\mu N=U$） \shoumonkai{2} $U=(AT/3)^{3/2}(NV)^{1/2}$，$p=(AT/3)^{3/2}(N/V)^{1/2}$，$F=-2(AT/3)^{3/2}(NV)^{1/2}$ \shoumonkai{3} $G=-N(AT/3)^{3}/p$ \shoumonkai{4} （確かめた）}
\begin{ronten}
  \item \shoumonkai{2} $F$ を $T$，$V$，$N$ の関数として書き直してから自然な変数の関数とみなすこと（$U$ や $S$ を残さないこと）を書く。
  \item \shoumonkai{3} $G$ の自然な変数は $T$，$p$，$N$ なので，状態方程式で $V$ を消してから答えることを書く。
\end{ronten}
\kenzan 示量性：$U$，$F$，$G$ はどれも $(V,N)\to(\lambda V,\lambda N)$ で $\lambda$ 倍になり，$p$，$\mu$ は変わらない。$\dd F=-S\,\dd T-p\,\dd V$ と整合（$-(\partial F/\partial T)_{V}=3U/T=S$）。
\tsumazuki $G=\mu N$ を確かめるとき，$\mu$ を $U,V,N$ のまま比べず，同じ変数 $(T,p)$ で書いてから比べる。
\end{tadashii}
\end{mondai}

\noteowari
\end{document}
